\documentclass[11pt,a4paper]{article}
\usepackage[T1]{fontenc}
\usepackage{lmodern}
\usepackage[margin=25mm]{geometry}
\usepackage{amsmath,amssymb,amsthm,mathtools}
\usepackage{booktabs,tabularx,array}
\usepackage{enumitem,microtype,xcolor}
\usepackage[colorlinks=true,linkcolor=blue!45!black,citecolor=blue!45!black,urlcolor=blue!45!black]{hyperref}
\usepackage[capitalise,nameinlink]{cleveref}
\usepackage{fancyhdr}
\pagestyle{fancy}
\fancyhf{}
\fancyhead[L]{\small Quantitative sieve exceptions and the Mills exponent}
\fancyhead[R]{\small KEIO}
\fancyfoot[C]{\thepage}
\setlength{\headheight}{14pt}
\newtheorem{theorem}{Theorem}[section]
\newtheorem{proposition}[theorem]{Proposition}
\newtheorem{lemma}[theorem]{Lemma}
\newtheorem{corollary}[theorem]{Corollary}
\theoremstyle{remark}
\newtheorem{remark}[theorem]{Remark}
\newcommand{\PP}{\mathcal P}
\newcommand{\W}{\mathcal W}
\newcommand{\meas}{\operatorname{meas}}
\newcommand{\Q}{\mathcal Q}
\newcommand{\LL}{\mathcal L}
\newcommand{\eps}{\varepsilon}
\newcommand{\N}{\mathbb N}
\newcommand{\R}{\mathbb R}
\newcommand{\Ch}{\operatorname{Ch}}
\title{\bfseries Lowering the Mills Exponent\\by Quantifying Sieve Exceptional Sets\\[6pt]
\large\normalfont An extension of harmonic branching to $c>43/41$}
\author{KEIO}
\date{8 November 2026}

\begin{document}
\maketitle
\begin{abstract}
Let $\W(c)$ consist of the real numbers $A>1$ for which
$\lfloor A^{c^n}\rfloor$ is prime for every $n\ge1$.
The harmonic branching manuscript supplied with this note obtains
$c>15/13$ from a short-interval prime number theorem.
We give a more economical analytic input: a positive proportion of the
expected primes, with sufficiently small exceptional measure.
We quantify the finite sieve argument of Li,
\emph{Primes in almost all short intervals}, arXiv:2407.05651v6,
with the free accuracy parameter
$\Q_D(X)=\exp(D(\log\log X)^2)$.
The Vinogradov--Korobov zero-free region supplies every fixed power of
$\Q_D^{-1}$ needed in the argument; all moment orders, sieve parameters,
and combinatorial depths are fixed before $X$ increases.
The resulting estimate, combined with harmonic branching, gives
$\W(c)\ne\varnothing$ for every $c>43/41=1.048780\ldots$,
and in fact uncountability and unboundedness.
The numerical region inequalities in Li's finite sieve decomposition are
an explicit external input, not independently recomputed in this note.
No Riemann hypothesis or density hypothesis is used.
The sharper potential targets $22/21$ and $24/23$ are recorded separately:
their required quantitative sieve estimates are not proved here.
Applying the same accuracy argument to Jia's published 1996 proof
already gives the independent range $c>20/19$.
\end{abstract}

\section{Result, scope, and dependencies}

Put
\[
 \W(c)=\{A>1:\lfloor A^{c^n}\rfloor\in\PP\text{ for every }n\ge1\},
 \qquad
 J_c(m)=[m^c,(m+1)^c-1).
\]
Here $c$ is the exponent parameter, and $A$ is the representing constant.
Neither the least element of $\W(3)$ nor its numerical value is changed.

\begin{theorem}[quantitative extension]\label{thm:main}
Using the finite sieve decomposition and numerical region bounds in
Li~\cite{LiI}, $\W(c)$ is uncountable, unbounded, Lebesgue-null and nowhere
dense for every real
\[
                       \boxed{c>\frac{43}{41}}.
\]
More generally, if $C_n=c_1\cdots c_n$ and
$43/41<\inf_n c_n\le\sup_n c_n<\infty$, the same assertions hold for
$\{A>1:\lfloor A^{C_n}\rfloor\in\PP\text{ for all }n\ge1\}$.
\end{theorem}

The qualification in the statement identifies the exact imported
component. Li's paper is a preprint, and its positive sieve margin rests
on explicit numerical integral inequalities. We use those inequalities
as stated; we do not supply a new interval-arithmetic certificate for
them. The new analytic work here concerns the uniform exceptional-set
estimate and its transfer to the prime tree. Thus this is a research
proof extension with documented dependencies, not a claim of independent
verification of every calculation in the cited literature.

\begin{theorem}[version using published sieve inputs]\label{thm:jia}
Using Jia's published finite sieve argument~\cite{Jia}, the conclusions
of \cref{thm:main} hold for $c>20/19$ and for bounded variable exponent
sequences whose infimum exceeds $20/19$.
This version does not require the improved regions or numerical
inequalities of Li's preprint.
\end{theorem}

\begin{center}
\small
\begin{tabularx}{\linewidth}{@{}lclX@{}}
\toprule
Input or target & interval exponent & Mills boundary & Status in this note\\
\midrule
Supplied manuscript & $2/15$ & $15/13$ & Baseline argument\\
Jia~\cite{Jia} & $1/20$ & $20/19$ & Quantitative extension of published sieve proof\\
Li I~\cite{LiI} & $2/43$ & $43/41$ & Quantitative extension proved below, using its finite sieve margin\\
Li II~\cite{LiII} & $1/22$ & $22/21$ & Further quantitative verification needed\\
Li III~\cite{LiIII} & $1/24$ & $24/23$ & Author-site manuscript; further quantitative verification needed\\
\bottomrule
\end{tabularx}
\end{center}

All displayed boundaries in this table concern strict inequalities.
In particular, $c=43/41$ is not included in \cref{thm:main}.
A convenient covered example is $c=21/20=1.05$, since
\[
 1-\frac1{21/20}=\frac1{21}>\frac2{43}.
\]
The reduction is
\[
 \frac{15}{13}-\frac{43}{41}=\frac{56}{533}=0.105065\ldots.
\]
Measured by the excess over $1$, it replaces $2/13$ by $2/41$.

\section{Why fixed logarithmic estimates cannot simply be substituted}

For each separately fixed $B$, Li's stated theorem supplies exceptional
measure $O_B(X(\log X)^{-B})$.
The harmonic construction needs a uniform estimate of the form
\[
             X\exp(-D(\log\log X)^2).
\]
The first assertion alone does not imply the second.
For example,
\[
             r(X)=\exp(-(\log\log X)^{3/2})
\]
is $O_B((\log X)^{-B})$ for every fixed $B$, but exceeds
$\exp(-D(\log\log X)^2)$ eventually for every fixed $D>0$.

This distinction has a concrete tree interpretation.
At level $n$, $\log Y_n$ is comparable to $c^n$, and repeatedly retaining
a prime proportion of size $1/\log Y_n$ loses
\[
 \prod_{j<n}\frac1{\log Y_j}
       =\exp\left(-\frac{\log c}{2}n^2+O(n)\right).
\]
A fixed logarithmic exceptional rate gives only
$\exp(-Bn\log c+O(1))$, which eventually exceeds this retained mass.
The analytic proof must therefore be made quantitative before it is
combined with the tree argument.

For $D>0$ define, for sufficiently large $X$,
\[
 \Q_D(X)=\exp(D(\log\log X)^2),\qquad
 \LL(X)=\frac{(\log X)^{1/3}}{(\log\log X)^{1/3}}.
\]
For any fixed $D,K,d,\nu>0$,
\begin{equation}\label{eq:hierarchy}
 (\log X)^d=\Q_D(X)^{o(1)},\quad
 \Q_D(X)^K=X^{o(1)},\quad
 \Q_D(X)^K e^{-\nu\LL(X)}=o(\Q_D(X)^{-d}).
\end{equation}
The order of choices below is essential: fix the interval exponent,
then fix all sieve and moment parameters, then choose the finite
accuracy powers, and finally let $X$ increase.

\section{A quantitative prime-polynomial estimate}

\begin{lemma}[uniform prime-polynomial bound]\label{lem:vk}
Fix $\delta>0$, $M_0>0$, and a bounded real interval of values of $\sigma$.
There is $a>0$ such that, whenever $X$ is sufficiently large,
$X^\delta\le u\le v\le2u\ll X$, and $|t|\le X^{M_0}$,
\[
 \sum_{u<n\le v}\frac{\Lambda(n)}{n^{\sigma+it}}
   =\int_u^v y^{-\sigma-it}\,dy
     +O_{\delta,M_0,\sigma}\left(u^{1-\sigma}e^{-a\LL(X)}\right).
\]
The constants are uniform when $\sigma$ stays in the indicated compact
interval. Consequently
\[
 \left|\sum_{u<n\le v}\frac{\Lambda(n)}{n^{\sigma+it}}\right|
 \ll u^{1-\sigma}
       \left(\frac1{1+|t|}+e^{-a\LL(X)}\right).
\]
The analogous upper bound for primes with coefficient $1$ holds, with at
most a fixed logarithmic factor.
\end{lemma}

\begin{proof}
Let $s=\sigma+it$. Apply truncated Perron's formula to
$-\zeta'(s+w)/\zeta(s+w)$ with difference kernel
\[
                    \frac{v^w-u^w}{w}.
\]
Take the initial line with $\Re(s+w)>1$, and choose a truncation height
$X^M$ with fixed $M$ sufficiently large relative to $M_0$ and the allowed
range of $\sigma$.
The endpoint convention for Perron's formula changes the normalized
sum by $O(u^{-\sigma}\log X)$, which is absorbed by the asserted error.

The Vinogradov--Korobov zero-free region permits a shift to
$\Re(s+w)=1-\eta$, where
\[
 \eta=b_0(\log X)^{-2/3}(\log\log X)^{-1/3}
\]
and $b_0>0$ is sufficiently small. Decreasing $b_0$ handles bounded
heights as well. In this narrower zero-free strip,
$-\zeta'/\zeta(z)-(z-1)^{-1}$ has a fixed power of $\log X$ as an upper
bound. The only pole crossed is at $w=1-s$, and its residue is
\[
              \frac{v^{1-s}-u^{1-s}}{1-s}=\int_u^v y^{-s}\,dy.
\]
The difference kernel is regular at $w=0$, so it creates no additional
residue there. The shifted vertical integral is
\[
 \ll u^{1-\sigma-\eta}(\log X)^{O(1)}
 \ll u^{1-\sigma}e^{-a\LL(X)},
\]
because $\log u\ge\delta\log X$. The horizontal integrals and Perron
truncation error are smaller on increasing $M$.
This proves the expansion.

For the main integral, writing $y=ue^z$ shows that its size is
$O(u^{1-\sigma}/(1+|t|))$, uniformly for $0\le\log(v/u)\le\log2$.
Removing prime powers costs at most a fixed logarithmic factor times
$u^{1/2-\sigma}+u^{-\sigma}$, which is negligible here.
Partial summation removes the remaining $\log p$ weight.
\end{proof}

\begin{corollary}\label{cor:Qinput}
For every fixed $D,K>0$, the prime polynomials in \cref{lem:vk} satisfy
\[
 \left|\sum_{u<n\le v}\Lambda(n)n^{-\sigma-it}\right|
 \ll u^{1-\sigma}\Q_D(X)^{-K}
\]
throughout $\Q_D(X)^{K+1}\le |t|\le X^{M_0}$, for sufficiently large $X$.
The low-frequency main integral is retained rather than discarded.
\end{corollary}

\begin{lemma}[fixed products of large primes]\label{lem:products}
The same conclusion, with a possible fixed logarithmic loss, holds for
Dirichlet polynomials formed from a fixed number of prime variables
$p_j\ge X^\delta$, with fixed order, range, or monomial constraints,
provided the selected prime variable has no arbitrary oscillating
coefficient.
It continues to hold after convolution with bounded sieve coefficients
in other variables, with fixed divisor and logarithmic losses.
\end{lemma}

\begin{proof}
Fix all but one prime variable and let $q$ be their product.
The allowed values of the remaining prime lie in a bounded union of
intervals inside $[M/q,2M/q]$, subject to a lower bound $X^\delta$.
Use \cref{lem:vk} on that variable. Its error, multiplied by $q^{-\sigma}$,
is $O(M^{1-\sigma}q^{-1}e^{-a\LL(X)})$.
The sum of $1/q$ over a fixed number of primes in $[X^\delta,O(X)]$ is
bounded in terms of $\delta$ and the number of variables, by Mertens'
estimate. Fixed logarithmic weights add only fixed logarithmic factors.

One can fix the other sieve variables first as well. Bounded sieve
coefficients have harmonic absolute sums bounded by fixed logarithmic
powers; grouping a fixed number of variables gives a fixed divisor
bound. These factors are absorbed using \eqref{eq:hierarchy}.
This statement does not assert cancellation for arbitrary coefficients
on the prime variable itself.
\end{proof}

\section{Quantifying the finite sieve argument}\label{sec:sieve}

We use Li's finite decomposition at the short-interval exponent $2/43$.
Its combinatorial regions and final numerical integrals are unchanged.
The issue is the analytic error budget.

\begin{proposition}[accuracy parameter in the finite sieve]\label{prop:finiteQ}
Fix a sufficiently small $\eps>0$ in the argument of \cite{LiI}, and put
\[
            \eta=\frac12X^{-41/43+\eps}.
\]
For every fixed $D>0$, there exist $X_0$ and $\kappa>0$ such that
\begin{equation}\label{eq:sieveQ}
 \meas\left\{x\in[X,2X]:
   \pi(x+\eta x)-\pi(x)<\kappa\frac{\eta x}{\log X}\right\}
 \le X\Q_D(X)^{-1}
 \qquad (X\ge X_0).
\end{equation}
Here the finite region inequalities (17), (18), and (20) of \cite{LiI}
are used as external numerical inputs.
\end{proposition}

\begin{proof}
Fix $\eps$ first; in particular $\delta=\eps^{1/3}$ is fixed.
Whenever necessary, take $\eps$ smaller so that the deterministic
sieve approximation errors fit strictly inside Li's positive margin.
No accuracy parameter below changes $\delta$, a moment order, or the
number of Buchstab iterations.

Write $Q=\Q_D(X)$.
We rerun the finite analytic estimates with an auxiliary accuracy
$Q^{-R}$, where $R$ is a sufficiently large fixed constant.
This is not an application of Li's theorem with a moving $B$.
The following checks show directly that its proof allows the new
accuracy.

\paragraph{1. Pointwise input.}
Replace the low-frequency cutoff by $T_0=Q^{R_0+2}$ with a sufficiently
large fixed $R_0$. Prime factors satisfy every fixed power of $Q^{-1}$
above that cutoff by \cref{cor:Qinput,lem:products}.
In particular, the saving $Q^{-R_0}$ remains available after the
fixed logarithmic losses.
Below it, keep the main integral in \cref{lem:vk}.
The contribution of its remainder is bounded by a fixed power of
$T_0$ and $\log X$ times $e^{-a\LL(X)}$, and is therefore smaller than
any prescribed fixed power of $Q^{-1}$.

\paragraph{2. Mean and large values.}
The estimates in Li's Lemmas 3.1--3.6 use fixed moments and fixed
interpolations of classical mean-value and large-value inequalities.
In Lemma 3.1 the integer moment satisfies $k\ll1/\delta$, and the
logarithmic loss has exponent $O(\delta^{-2})$; both are fixed.
Every resulting case has one of two forms:
\[
     (\log X)^d Q^{-\alpha R_0}
             \quad\hbox{or}\quad
     (\log X)^d X^{-\nu},
\]
where $d$ is fixed, $\alpha>0$, and $\nu>0$ depends on the fixed
exponent margin. There are finitely many such cases.
Thus the minimum relevant positive $\alpha$ can be chosen before
$R_0$. By increasing $R_0$ and then $X$, every one is
$O(Q^{-R})$ for any specified fixed $R$.
The extra finite polynomial products in Lemmas 3.3 and 3.4 do not
alter this conclusion. Lemmas 3.5 and 3.6 are regroupings and
permutations of Lemma 3.2 and preserve the same accuracy accounting.

For completeness, one displayed interpolation in Lemma 3.2,
case (1)(a), can be justified without its copied intermediate line.
If the simultaneous large-value count $F$ is bounded by
\[
 V^{-6}M^3,\quad W^{-10}H^5,\quad U^{-4}K^2,
 \qquad MHK=X,
\]
interpolate with weights $3/10$, $1/5$, and $1/2$. Then
\begin{align*}
 U^2V^2W^2F
 &\ll U^2V^2W^2
     (V^{-6}M^3)^{3/10}
     (W^{-10}H^5)^{1/5}(U^{-4}K^2)^{1/2}\\
 &=V^{1/5}M^{9/10}HK.
\end{align*}
If $V\le M^{1/2}Q^{-R_0}$, this is at most
$XQ^{-R_0/5}$. The positive power of the pointwise saving is all that
is needed. This is the interpolation already present in Jia's
Lemma 5, case I(a), printed page 28.

\paragraph{3. The Watt-type estimates.}
Li's Section 4 uses a fixed exponent margin in its mean-value
estimates. The resulting error is a fixed power saving in $X$,
or one of the small-factor cases already covered in step 2.
Both dominate every required fixed power of $Q^{-1}$ by
\eqref{eq:hierarchy}. The constants in these mean-value estimates
depend on the fixed margins and coefficient bounds, not on a
logarithmic accuracy exponent.

\paragraph{4. Perron and the exceptional measure.}
Here are the details behind the required error bookkeeping.
For a relevant Dirichlet polynomial $F$, put
\[
                W_\eta(s)=\frac{(1+\eta)^s-1}{s}.
\]
On a dyadic range $|t|\asymp T$,
$|W_\eta(b+it)|\ll\min(\eta,T^{-1})$, where $b=1+1/\log X$.
Opening the square, integrating $x^{s+\overline{s'}}$, and applying
$2|ab|\le |a|^2+|b|^2$, gives
\begin{align}
 &\int_X^{2X}\left|
     \int_{T_0\le |t|\le X}F(b+it)W_\eta(b+it)x^{b+it}\,dt
                       \right|^2dx \nonumber\\
 &\hspace{7mm}\ll X^3(\log X)^3
       \max_{T_0\le T\le X}
       \min(\eta,T^{-1})^2
          \int_{T\le |t|\le2T}|F(b+it)|^2dt. \label{eq:perronL2}
\end{align}
This is the estimate used in Jia's Lemmas 16--19~\cite{Jia}, whose
Perron steps underlie Li's Section 5.
If the right-hand maximum is $\ll\eta^2Q^{-R}$, Chebyshev's inequality,
with threshold $\eta XQ^{-r}$, yields exceptional measure
\[
                       \ll X(\log X)^3Q^{-R+2r}.
\]
Choose $R>2r+10$ after all fixed decompositions. This leaves more
than enough accuracy for \eqref{eq:sieveQ}.
The low-frequency terms are evaluated as main terms, with errors
$\ll\eta XQ^{-r}$ by step 1.
Terms such as $\eta^2XT_0^j$ and Perron truncation errors are harmless
since $T_0^j=X^{o(1)}$ and $\eta X$ is a fixed positive power of $X$.
Endpoint conventions contribute smaller errors as well.
The same argument with the appropriate zeta factor handles the
Type I terms.

\paragraph{5. Separation of the finite constraints.}
The number of prime variables and the number of defining inequalities
are fixed. Standard dyadic partitions introduce only a fixed power
of $\log X$. Alternatively, a relative mesh $Q^{-r_1}$ separates the
variables into at most $Q^{O_\eps(r_1)}$ boxes.
Boundary strips in fixed monomial and ordering constraints have
total tuple weight $\ll XQ^{-r_1}(\log X)^{O_\eps(1)}$; fixing all but
one variable proves this bound by the elementary divisor estimate.
Fubini and Markov transfer it to the short sums.
Choose $r_1$ for the needed boundary accuracy first, and then choose
$R$ large enough to cover the number of boxes.
This avoids a circular choice of accuracy powers.
For the lower bound one may also shrink the admissible polyhedral
regions by a sufficiently small fixed amount: their boundary loss is
absorbed in the fixed positive numerical margin.

\paragraph{6. The sieve and its positive margin.}
In Li's Lemma 5.6 the sieve level $X^{10^{-1000}}$ and the sifting
threshold $X^\delta$ are fixed powers of $X$.
The bounded sieve coefficients and the grouping of finitely many
variables preserve fixed divisor bounds. The deterministic errors of
size $O(\eps)$ in the normalized main term need not become $Q^{-1}$:
choose $\eps$ small once. Only the fluctuating error and its exceptional
set require the stronger accuracy.

In applying the small-factor estimates, a coefficient supported on
rough integers is not automatically declared to have cancellation.
The relevant groups contain prescribed prime factors, to which
\cref{lem:products} applies after the remaining variables are fixed.
If the group denoted $N$ in Li's Lemma 4.4 consists only of the tiny
sieve variable, use the Type I estimate in Lemma 5.1 instead:
the relevant range has $QR<X^{45/86}$ and $N<X^{41/344}$.
Thus the proof does not need a pointwise estimate for arbitrary
sieve weights.

Li's recorded numerical region estimates give the positive margin
\[
                  1-0.03-0.992+0.027=0.005.
\]
Choose the fixed deterministic errors and region-boundary losses less
than half this margin. Steps 1--5 make all remaining errors smaller
than a further fixed portion outside measure $XQ^{-1}$, for sufficiently
large $X$. The same fixed sieve decomposition therefore gives
\eqref{eq:sieveQ}, with a positive constant $\kappa$.
The numerical inequalities themselves are precisely the external
component specified in the proposition.
\end{proof}

\begin{remark}[what is and is not imported]
The standard statement $O_B(X/\log^B X)$ is not the imported
quantitative conclusion. We use the finite algebraic estimates and
the fixed positive sieve margin, and verify their accuracy dependence.
The general free-parameter variance estimate in
Matomäki--Radziwiłł--Shao--Tao--Teräväinen~\cite[Lemma 3.3]{MRSTT}
provides a separate modern formulation of the Perron/variance step.
Their other results are not used to claim a new short-interval exponent.
\end{remark}

\section{Transfer from short intervals to bad integer nodes}

\begin{proposition}[the published-source alternative]\label{prop:jiaQ}
For sufficiently small fixed $\eps>0$, the counterpart of
\cref{prop:finiteQ} holds with
$\eta=\tfrac12X^{-19/20+\eps}$, using Jia's published proof
in place of Li's refined decomposition.
\end{proposition}

\begin{proof}
Jia fixes $\eps$, $\eps_1=\eps^2$ and $\delta=\eps^{1/3}$ before
the accuracy parameter. His mean-value arguments in Sections 2--5
have the same structure as steps 1--3 in \cref{prop:finiteQ}:
fixed moments, fixed powers of logarithms, and either a positive
power of the pointwise saving or a fixed power saving in $X$.
His Lemmas 16--19 are the Perron calculations quantified directly
in \eqref{eq:perronL2}. The subsequent fixed Buchstab regions,
linear sieve levels and deterministic $O(\eps)$ or $O(\delta)$
errors do not depend on the requested logarithmic accuracy.
Thus the identical order of parameter choices gives an exceptional
measure at most $X\Q_D(X)^{-1}$ for every fixed $D$.
Jia's published positive numerical sieve margin supplies the
constant $\kappa>0$. We use that published margin as an external
input just as in a normal application of his theorem.
\end{proof}

A separate averaging argument avoids any need to force the sieve
interval length to equal the length of $J_c(m)$.

\begin{lemma}[longer interval averaging]\label{lem:average}
Suppose, for $u\in[X,2X]$ outside a set $E_X$,
\[
       \#\bigl(\PP\cap(u,(1+\eta)u]\bigr)
                       \ge\kappa\frac{\eta u}{\log X}.
\]
Let $I=[a,b)\subset[X,2X]$, $L=b-a$, and $L/(\eta X)\to\infty$.
If $\#(\PP\cap I)<\kappa L/(16\log X)$, then
\[
                         \meas(E_X\cap I)\ge L/4
\]
for sufficiently large $X$.
\end{lemma}

\begin{proof}
Let $B_I=[a,b/(1+\eta))$.
For large $X$, $\meas B_I\ge L/2$ and
$(u,(1+\eta)u]\subset I$ for $u\in B_I$.
Every prime of $I$ is counted for at most $2\eta X$ units of $u$-measure.
Consequently
\[
 \int_{B_I}\#\bigl(\PP\cap(u,(1+\eta)u]\bigr)\,du
 \le2\eta X\#(\PP\cap I)<\frac{\kappa\eta XL}{8\log X}.
\]
On $B_I\setminus E_X$ the integrand is at least
$\kappa\eta X/\log X$, so
$\meas(B_I\setminus E_X)\le L/8$.
Thus $\meas(E_X\cap I)\ge3L/8\ge L/4$.
\end{proof}

\begin{proposition}[bad-node estimate]\label{prop:bad}
Fix $43/41<c_-\le c_+<\infty$.
There is $\delta_0>0$ such that, for every $D>0$, for sufficiently
large $Y$, uniformly for $c\in[c_-,c_+]$,
\[
 \#\left\{m\in[Y,2Y):
 \#(\PP\cap J_c(m))<
     \delta_0\frac{|J_c(m)|}{\log((m+1)^c)}\right\}
 \le Y\exp(-D(\log\log Y)^2).
\]
\end{proposition}

\begin{proof}
Choose a fixed sufficiently small $\eps>0$ with
\[
            \frac2{43}+\eps<1-\frac1{c_-}.
\]
At height $X\asymp Y^c$, the intervals $J_c(m)$ have length
\[
            |J_c(m)|\asymp Y^{c-1}=X^{1-1/c}.
\]
This is uniformly larger by a fixed power of $X$ than the intervals
of length $\eta X\asymp X^{2/43+\eps}$ in \cref{prop:finiteQ}.
The $J_c(m)$ are disjoint.
Cover their heights with $O_{c_+}(1)$ dyadic ranges.
At most $O_{c_+}(1)$ of the intervals cross a boundary; absorb these
in the final estimate.

Choose $\delta_0$ sufficiently small in terms of $\kappa$.
By \cref{lem:average}, each bad $m$ accounts for
$\gg Y^{c-1}$ measure of the exceptional set in its dyadic range.
Apply \cref{prop:finiteQ} with a larger fixed accuracy constant, say
$2D+1$, and divide the exceptional measure by $Y^{c-1}$.
Since
$\log\log(Y^c)=\log\log Y+\log c$, fixed factors and the dyadic cover
are absorbed for large $Y$. This proves the estimate.
\end{proof}

\section{The harmonic construction}

We include the mass argument to specify the required quantifiers.

\begin{lemma}[mass survival and binary subdivision]\label{lem:mass}
Let $V_1$ be finite, and let disjoint finite child sets $\Ch_n(m)$
define the successive levels
$V_{n+1}=\bigcup_{m\in V_n\cap G_n}\Ch_n(m)$.
Write $w(S)=\sum_{m\in S}1/m$.
Suppose
\[
 w(\Ch_n(m))\ge\lambda_n/m\quad(m\in V_n\cap G_n),
 \qquad w(V_n\setminus G_n)\le\beta_n.
\]
Put $s_n=\frac13w(V_1)\prod_{j<n}(\lambda_j/6)$.
If $\beta_n\le s_n/2$ and $\max_{m\in V_n}1/m\le s_n$ for all $n$,
there are uncountably many infinite branches.
\end{lemma}

\begin{proof}
A finite collection of total weight at least $3s_n$, with every
atom of weight at most $s_n$, splits into two disjoint parts of
weight at least $s_n$ each, by a greedy cut.
After selecting either part and discarding its bad nodes, its
remaining weight is at least $s_n/2$.
Its children have weight at least
$\lambda_ns_n/2=3s_{n+1}$, so the operation repeats.
For every infinite binary choice the selected forest has nodes at
every depth. Local finiteness and the finite root set give an infinite
path by König's lemma.
Different binary choices have disjoint selected node sets at the
first differing level, hence give different paths.
\end{proof}

\begin{proof}[Proof of \cref{thm:main}, constant exponent]
Fix $c>43/41$ and a large $P$.
Start from $V_1=\PP\cap[P,2P]$ and put
$\Ch_n(m)=\PP\cap J_c(m)$.
Take $G_n$ to be the good integers in \cref{prop:bad}.
Inductively,
\[
             V_n\subset[P^{c^{n-1}},(3P)^{c^{n-1}}]
\]
for large $P$. The upper bound follows by iterating
$M_{n+1}=(M_n+1)^c$ and summing the increments $1/M_n$ after
dividing logarithms by $c^{n-1}$.

The prime number theorem gives
$w(V_1)\ge1/(4\log P)$.
Since $|J_c(m)|\sim cm^{c-1}$, goodness gives
\[
             w(\Ch_n(m))
      \ge\frac{a}{m\,c^{n-1}\log(4P)}
\]
with fixed $a>0$.
Set $\lambda_n=a/(c^{n-1}\log(4P))$.
A dyadic cover of the displayed level range has $O(c^{n-1})$
pieces, so \cref{prop:bad} yields
\[
 \beta_n\ll c^{n-1}
       \exp\left(-D\bigl(\log\log(P^{c^{n-1}})\bigr)^2\right).
\]

Choose $D>1/(2\log c)$.
Write $m=n-1$, $d=\log c$, $\ell=\log P$, and $L=\log\ell$.
For $s_n$ in \cref{lem:mass},
\[
 -\log s_n=\frac d2m^2+mL+O(m+L+1),
\]
whereas
\[
 \log\beta_n\le-D(L+md)^2+md+O(1).
\]
Their difference is at most
\[
 -\left(Dd^2-\frac d2\right)m^2
 -(2Dd-1)mL-DL^2+O(m+L+1).
\]
Both bracketed coefficients are positive. A single sufficiently
large $P$ therefore ensures $\beta_n\le s_n/2$ for every $n$.
Also
\[
       \max_{p\in V_n}\frac1p\le e^{-c^m\ell}\le s_n
\]
for all $m\ge0$ once $P$ is sufficiently large.
Thus \cref{lem:mass} supplies uncountably many infinite prime chains.

For such a chain, set
\[
       a_n=p_n^{1/c^n},\qquad b_n=(p_n+1)^{1/c^n}.
\]
The inequalities defining $J_c(p_n)$ imply
$a_{n+1}\ge a_n$ and $b_{n+1}<b_n$.
Hence $A=\lim a_n$ exists and
$a_n\le A\le b_{n+1}<b_n$.
Thus $\lfloor A^{c^n}\rfloor=p_n$ for every $n$.
Different chains yield different $A$.
Since $P$ can be arbitrarily large, $\W(c)$ is unbounded.

Finally, fix $1<a<b$. The set satisfying just the $n$th prime
condition in $[a,b]$ is a finite union of intervals of total length
at most
\[
       \frac b{c^n}
          \sum_{a^{c^n}-1\le p\le b^{c^n}}\frac1p=O_{a,b}(c^{-n}),
\]
by Mertens' theorem.
Its closure adds only finitely many endpoints.
The closure of $\W(c)\cap[a,b]$ is contained in every such closed
union, hence is null and has empty interior.
This proves nullity and nowhere density.
\end{proof}

\begin{proof}[Bounded variable exponents]
Let $D_n=c_2\cdots c_n$, $D_1=1$.
Replace $c^{n-1}$ above by $D_n$.
The uniform estimate in \cref{prop:bad} applies because
$c_n\in[c_-,c_+]$.
Use $\lambda_n=a/(D_n\log(4P))$ and choose
\[
                 D>\frac{\log c_+}{2(\log c_-)^2}.
\]
The same comparison, using $c_-^{n-1}\le D_n\le c_+^{n-1}$, proves
the conditions of \cref{lem:mass} uniformly at every level.
The nested intervals now have exponents $1/C_n$.
Their limits, distinctness, unboundedness, and the nullity argument
are unchanged.
\end{proof}

\begin{proof}[Proof of \cref{thm:jia}]
Use \cref{prop:jiaQ} in the sliding-window averaging argument.
The only change in the bad-node and harmonic-tree proofs is the
choice $1/20+\eps<1-1/c_-$, which is possible exactly when
$c_->20/19$. All other estimates are the same.
\end{proof}

\section{What would lower the exponent further?}

The extension above uses the first Li paper because its finite
mean-value calculation and the Perron arguments behind it can be
tracked with an explicit accuracy parameter.
Li II states a lower-bound theorem at interval exponent $1/22$.
Its additional ingredients include mean-value estimates quoted from
Harman's monograph; the author-site manuscript Li III gives the
further proposed exponent $1/24$.
Both state fixed-$B$ exceptional estimates.

The exact next task is to prove, for their additional sieve regions,
the counterpart of \cref{prop:finiteQ}:
\[
 \forall D>0,\qquad
 \meas\{x:\text{the required lower bound fails}\}
             \ll X\exp(-D(\log\log X)^2).
\]
If established at every exponent $\theta>1/22$, it gives
$c>22/21$; at every $\theta>1/24$, it gives $c>24/23$.
These are implications, not conclusions of the present proof.

The following distinctions should be preserved in any subsequent
revision:
\begin{itemize}[leftmargin=1.5em]
\item A fixed finite prime chain is not evidence of an infinite chain.
\item A new exceptional-set estimate must be proved with its parameter
      dependence; choosing $B$ as a function of $X$ in a fixed-$B$
      theorem is insufficient.
\item The representing constant obtained by compactness need not be
      the limit of a specified greedy prime search.
\item The endpoint $c=43/41$ and the whole range $1<c\le43/41$
      remain outside the proved conclusion of this note.
\item No historical-priority or publication claim is made.
\end{itemize}

\section*{Dependency record}

\begin{enumerate}[leftmargin=1.8em]
\item The supplied harmonic branching manuscript~\cite{KEIO}
      provides the original global-mass construction.
      Its core survival calculation is reproduced here.
\item Li I~\cite{LiI} supplies the finite admissible sieve regions,
      the large-value case structure, and the numerical integral
      bounds giving the positive margin $0.005$.
\item Jia~\cite{Jia} supplies the underlying Perron and variance
      computations; their accuracy dependence is made explicit in
      \eqref{eq:perronL2}.
\item Vinogradov--Korobov supplies the stronger pointwise error used
      in \cref{lem:vk}. All other losses in the audited finite
      argument are fixed powers or fixed positive powers of the
      new accuracy target.
\item The more recent manuscript thresholds are documented as
      research targets, not silently imported as already quantified
      theorems.
\end{enumerate}

\begin{thebibliography}{99}\small\raggedright
\bibitem{KEIO}
KEIO, \emph{Harmonic Branching Mass and Mills-Type Prime-Representing
Functions below the Exponent $20/11$}, supplied manuscript,
8 November 2026, file \texttt{harmonic\_branching\_mills.tex}.
\bibitem{LiI}
R. Li, \emph{Primes in almost all short intervals},
arXiv:2407.05651v6, 25 September 2025.
\url{https://arxiv.org/abs/2407.05651v6}.
\bibitem{LiII}
R. Li, \emph{Primes in almost all short intervals II},
Cambridge Open Engage preprint, 2025.
\url{https://doi.org/10.33774/coe-2025-jrnjl}.
\bibitem{LiIII}
R. Li, \emph{Primes in almost all short intervals III},
author-site manuscript, accessed 8 November 2026.
\url{https://runbolicarey.com/assets/downloads/Primes_in_almost_all_short_intervals_III.pdf}.
\bibitem{Jia}
C. Jia, Almost all short intervals containing prime numbers,
\emph{Acta Arithmetica} 76 (1996), 21--84.
\url{https://matwbn.icm.edu.pl/ksiazki/aa/aa76/aa7613.pdf}.
\bibitem{MRSTT}
K. Matomäki, M. Radziwiłł, X. Shao, T. Tao, J. Teräväinen,
Higher uniformity of arithmetic functions in short intervals II.
Almost all intervals, \emph{Inventiones Mathematicae} 244 (2026),
967--1091.
\url{https://doi.org/10.1007/s00222-026-01408-6}.
\bibitem{Harman}
G. Harman, \emph{Prime-Detecting Sieves},
London Mathematical Society Monographs 33,
Princeton University Press, 2007.
\bibitem{Ford}
K. Ford, Vinogradov's integral and bounds for the Riemann zeta
function, \emph{Proc. London Math. Soc.} (3) 85 (2002), 565--633.
\bibitem{GM}
L. Guth, J. Maynard, New large value estimates for Dirichlet
polynomials, \emph{Annals of Mathematics} 203 (2026), 623--675.
\url{https://arxiv.org/abs/2405.20552}.
\end{thebibliography}
\end{document}

